% Performance-analysis section for the Qsim report.
% Requires only the article class (no extra packages).
% Raw data: verify/bench.txt, produced by `go run ./cmd/bench`.

\section{Performance Analysis}
\label{sec:performance-analysis}

To characterize the performance limits of QSim, the computational cost of
the simulation engine was measured headless on an 11th Gen Intel Core
i7-1165G7 processor with 16\,GB of memory; reported values are best-of-$N$
runs. The benchmark exercises the same primitives the desktop application
executes on every frame when recomputing gate outputs -- state merging,
qubit reordering, and unitary application on complex64 state vectors -- so
the pipeline times below constitute a lower bound on per-frame engine
cost, excluding rendering and interface work.

Gate latency (Table~\ref{tab:gate-latency}) is the cost of applying one
dense $k$-qubit gate to an $n$-qubit state. Cost grows by roughly a factor
of four per additional state qubit and per additional gate qubit, as
expected for dense unitary application: up to ten qubits every gate width
completes below timer resolution, while at sixteen qubits a single-qubit
gate costs about $1$\,ms and a four-qubit universal gate about $12$\,ms.

\begin{table}[h]
\centering
\caption{Dense $k$-qubit gate latency (ms) on an $n$-qubit state.}
\label{tab:gate-latency}
\begin{tabular}{c|c|cccc}
$n$ & mem & $k=1$ & $k=2$ & $k=3$ & $k=4$ \\ \hline
$\le 10$ & $<0.1$\,MB & $<0.01$ & $<0.01$ & $<0.01$ & $<0.01$ \\
12 & $<0.1$\,MB & $<0.01$ & $<0.01$ & $<0.01$ & 0.53 \\
14 & 0.1\,MB & $<0.01$ & $<0.01$ & 1.00 & 3.00 \\
16 & 0.5\,MB & 1.00 & 2.99 & 5.99 & 12.00 \\
\end{tabular}
\end{table}

For a representative frame of engine work -- merging the single-qubit
systems, one Hadamard layer, and one CNOT ladder -- a full recomputation
(Table~\ref{tab:pipeline}) costs far below a millisecond up to ten qubits,
$4.00$\,ms at twelve qubits, $15.59$\,ms at fourteen, and $80.00$\,ms at
sixteen, with cost per added qubit reaching four to five times beyond that
point. Against a 16.7\,ms frame budget the qubit ceiling is therefore
approximately fourteen qubits at 60\,fps, and approximately fifteen qubits
at 30\,fps.

\begin{table}[h]
\centering
\caption{Representative-frame engine cost (ms) vs qubit count.}
\label{tab:pipeline}
\begin{tabular}{c|c|ccc|c|c}
$n$ & mem & merge & H layer & CX ladder & total & $\approx$fps \\ \hline
$\le 10$ & $<0.1$\,MB & $\approx 0$ & $\approx 0$ & $\approx 0$ & $\approx 0$ & $>10000$ \\
12 & $<0.1$\,MB & 0.00 & 2.00 & 2.00 & 4.00 & 250 \\
14 & 0.1\,MB & 0.00 & 9.59 & 6.01 & 15.59 & 64 \\
16 & 0.5\,MB & 1.00 & 46.00 & 33.00 & 80.00 & 13 \\
18 & 2.1\,MB & 4.00 & 188.61 & 141.09 & 333.70 & 3 \\
20 & 8.4\,MB & 17.11 & 909.12 & 626.71 & 1552.94 & 1 \\
\end{tabular}
\end{table}

These ceilings are engine-only: in-app rendering draws one cell per
amplitude, so the real interactive ceiling lies at or below the reported
values. Memory is never the binding constraint -- the state vector
occupies only 8.4\,MB at twenty qubits while compute already exceeds one
and a half seconds -- and individual frames vary with gate count and
width. Agreement between the measured scaling and the theoretical
expectations for state-vector simulation increases confidence that no
hidden overhead dominates the engine.
